\documentclass{handout}
\usepackage{handoutSetup}\usepackage{handoutShortcuts}  

%\setboolean{BigAndWide}{true}\wideMargins  % Useful for making a big-print version that can be read easily even when printed in a two-pages-per-paper-side format

\opt{bigWide}{\wideMargins}
\begin{document}
\handoutHeader


\begin{verbatimwrite}{\jobname.title}
Generational Accounts and the Government
\end{verbatimwrite}

\handoutNameMake


\section{The Government Budget Constraint}

Consider a government that raises taxes $\TaxLev_{t}$, makes expenditures 
$\GovSpend_{t}$, and has an outstanding stock of debt $\Debt_{t}$ at the beginning 
of period $t$, on which it must pay interest at rate $\rfree_{t}$.  The 
government can run a deficit only by raising funds via the issuing of 
new bonds.\opt{MarginNotes}{\marginpar{\tiny Other way would be to print money - see Larry Ball's class for this.}}

The government's Dynamic Budget Constraint (DBC) is given by 
\begin{equation}\begin{gathered}\begin{aligned}
        \overbrace{\Debt_{t+1}-\Debt_{t}}^{\mbox{Deficit}} & =  \overbrace{(\GovSpend_{t}+\rfree_{t}\Debt_{t})}^{\mbox{Outlays}}-\TaxLev_{t} % \label{eq:govdbc}
\\ \Debt_{t+1} & =  \GovSpend_{t}+\Rfree_{t}\Debt_{t}-\TaxLev_{t}
\\ \Debt_{t} & =  \left(\frac{\Debt_{t+1}+\TaxLev_{t}-\GovSpend_{t}}{\Rfree_{t}}\right)   \label{eq:recurse}
.
\end{aligned}\end{gathered}\end{equation}

But we can obtain a similar formula for $\Debt_{t+1}$ in terms of $\Debt_{t+2}$,
and substitute it into \eqref{eq:recurse}.  Continued substitution
gives\opt{MarginNotes}{\marginpar{\tiny This equation just says gov must plan to repay its debts.  No constraint on what gov can do in any one or 2 periods.  Diff from HHs in this way.  Because gov has $\infty$ horizon, can obligate future generations to pay for spending of current generations.  Quite a trick!  (Japan 100 year mortgage).}}
\begin{equation}\begin{gathered}\begin{aligned}
        \Debt_{t} & =  \overbrace{(\TaxLev_{t}-\GovSpend_{t})}^{\equiv \Surplus_{t}}/\Rfree_{t}+(\TaxLev_{t+1}-\GovSpend_{t+1})/\Rfree_{t}\Rfree_{t+1}+\ldots  \\
              & =  \Surplus_{t}/\Rfree_{t}+\Surplus_{t+1}/\Rfree_{t}\Rfree_{t+1}+\ldots  \\ 
        \Rfree_{t}\Debt_{t} & =  \mathbb{P}_{t}(\Surplus)   \\
         & =  \mathbb{P}(\mbox{Govt Primary Surpluses}) \label{eq:govibc}
\end{aligned}\end{gathered}\end{equation}
where $\mathbb{P}$ denotes the present discounted value; this can be rewritten
\begin{equation}\begin{gathered}\begin{aligned}
        \mathbb{P}_{t}(X) & =  \mathbb{P}_{t}(T)-\Rfree_{t}\Debt_{t} \label{eq:govibc2}.
\end{aligned}\end{gathered}\end{equation}

Equation \eqref{eq:govibc2} should look familiar: recall that in the consumption problem we 
had an Intertemporal Budget Constraint that said  \opt{MarginNotes}{\marginpar{\tiny $\mathbb{P}_{t}(Y) = H_{t}$.}}
\begin{equation}\begin{gathered}\begin{aligned}
        \mathbb{P}_{t}(C) & =  \mathbb{P}_{t}(Y)+\Rfree_{t} \Kap_{t}
\end{aligned}\end{gathered}\end{equation}
where $\Kap_{t}$ is the beginning-of-period level of capital wealth (before
interest has been earned).  

\opt{MarginNotes}{\marginpar{\tiny It's $RB$ rather than $B$ 
because we set up budget constraint with timing of interest payments 
at beginning of period rather than end.  Note that as time interval goes
to zero, $R \approx 1$.
}}
In each case, the PDV of expenditures must be equal to the PDV of 
income plus current wealth.  Thus, equations \eqref{eq:govibc} and 
\eqref{eq:govibc2} are different ways to express the Government 
Intertemporal Budget Constraint (GIBC).\footnote{The $\Rfree_{t}$ is present
in the government's problem because for the consumer we were thinking
about the situation after any interest income was received; if we
were to think of the consumer's beginning-of-period capital as $K_{t}$
then we would have $B_{t} = \Rfree_{t} K_{t}$; and note that the sign difference
reflects the fact that $D$ is debt while $B$ is balances.}

Now let's suppose that the only kind of expenditures the government  \opt{MarginNotes}{\marginpar{\tiny This is net transfers in year $t$ (if we are only thinking of SS), `primary surplus' if we think of govt activities more generally.}} 
engages in are transfers, so that $\GovSpend_{t}$ simply reflects money handed 
out to some members of the population in period $t$.  Then $\Surplus_{t}$ 
will be equal to total net transfers among the members of the 
population at period $t$.  Note 
that there is nothing that says that $\Surplus_{t}$ must be positive or 
negative in any particular period.  The GIBC only places restrictions 
on the present discounted value of net transfers.

The fact that government only has to satisfy the GIBC means that the 
government can potentially treat different generations very 
differently from each other.  It is therefore useful to have a 
mechanism to keep track of how different generations are treated.  The 
standard way of doing this is to construct a set of `generational 
accounts,' as initially proposed by \cite{akg:genaccts}.

If we assume that consumers live two-period lives, the generational
account for the generation born at time $t$ is:\opt{MarginNotes}{\marginpar{\tiny Point out that either of these $Z$'s can be negative, which would signify net transfers to households.}}
\begin{equation}\begin{gathered}\begin{aligned}
        \bar{Z}_{t} & =  \Surplus_{1,t}+\Surplus_{2,t+1}/\Rfree_{t+1}  \\
         & =  \mbox{PDV of lifetime taxes net of transfer payments}.
\end{aligned}\end{gathered}\end{equation}

In the US and most other countries, working-age people pay more in 
taxes than they receive in transfers, so $\Surplus_{1,t}$ is positive, while 
old people receive more in transfers than they pay in taxes, so 
$\Surplus_{2,t}$ is negative.\opt{MarginNotes}{\marginpar{\tiny Largest transfers are Social Security and Medicare in US.}}

Note now that the aggregate total of net transfers can be subdivided 
into the net transfers of the two age groups in the population,
\begin{equation}\begin{gathered}\begin{aligned}
        \Surplus_{t} & =  \Surplus_{1,t}+\Surplus_{2,t}.
\end{aligned}\end{gathered}\end{equation}
%and note further that the GIBC can be rewritten as
%\begin{equation}\begin{gathered}\begin{aligned}
%        \Rfree_{t}\Debt_{t} & =  \mathbb{P}_{t}(T) - \mathbb{P}_{t}(X) = \mathbb{P}_{t}(T-X) = \mathbb{P}_{t}(Z)
%\end{aligned}\end{gathered}\end{equation}

Now write out the GIBC \eqref{eq:govibc} explicitly:
\begin{equation}\begin{gathered}\begin{aligned}
        \mathbb{P}_{t}(Z) & =  \phantom{+} \Surplus_{t}+\Surplus_{t+1}/\Rfree_{t+1}+ \ldots \\
         & =  \phantom{+} \Surplus_{1,t}+\Surplus_{1,t+1}/\Rfree_{t+1}+\Surplus_{1,t+2}/\Rfree_{t+1}\Rfree_{t+2}+\ldots \\
         &    +\Surplus_{2,t}+\Surplus_{2,t+1}/\Rfree_{t+1}+\Surplus_{2,t+2}/\Rfree_{t+1}\Rfree_{t+2}+\ldots 
         \nonumber \\
         & =  \Surplus_{2,t}+[\Surplus_{1,t}+\Surplus_{2,t+1}/\Rfree_{t+1}]+[\Surplus_{1,t+1}+\Surplus_{2,t+2}/\Rfree_{t+2}]/\Rfree_{t+1} + \ldots
\\       & =  \Surplus_{2,t}+\bar{Z}_{t}+\bar{Z}_{t+1}/\Rfree_{t+1}+\bar{Z}_{t+2}/\Rfree_{t+1}\Rfree_{t+2}+\ldots
\end{aligned}\end{gathered}\end{equation}
which again shows that the GIBC is consistent with any treatment of 
any particular generation; any pattern of generational accounts that 
satisfies the GIBC is feasible.

\section{Social Security and Generational Accounts}

Consider an economy that initially has no government so that 
$\Surplus_{1,t}=\Surplus_{2,t}=\Surplus_{2,t-1} = 0$.  Now consider introducing a Pay As 
You Go (PAYG) Social Security system at date $s$, which is to remain 
of constant size forever after introduction,
\begin{equation}\begin{gathered}\begin{aligned}
        \Surplus_{2,t} & =   -\Surplus_{1,t} \neq 0~~ \forall~t~\geq~s
\\      \Surplus_{1,t+1} & =  \Surplus_{1,t}.
\end{aligned}\end{gathered}\end{equation}

Consider the generation born at time $s-1$.  It paid nothing into the 
Social Security system when young, yet gets $\Surplus_{2,s}$ out when old.
Its generational account is therefore 
\begin{equation}\begin{gathered}\begin{aligned}
        \bar{Z}_{s-1} & =  \Surplus_{1,s-1}+\Surplus_{2,s}/\Rfree_{s}  \\
         & =  0+\Surplus_{2,s}/\Rfree_{s}
\end{aligned}\end{gathered}\end{equation}
so this generation benefits from the introduction of SS because \opt{MarginNotes}{\marginpar{\tiny Since $\Surplus_{2,s}$ is a neg number, they are better off than without SS.}}
it paid no taxes yet receives benefits.

The GA's for succeeding generations are
\begin{equation}\begin{gathered}\begin{aligned}
        \bar{Z}_{t} & =  \Surplus_{1,t}+\Surplus_{2,t+1}/\Rfree_{t+1}  \\
         & =  \Surplus_{1,t}(1-1/\Rfree_{t+1})  \\
         & =  \rfree_{t+1}\Surplus_{1,t}/\Rfree_{t+1} \label{eq:SScost}
\end{aligned}\end{gathered}\end{equation}
so future generations are worse off by this amount.\opt{MarginNotes}{\marginpar{\tiny Taxes are positive; effect on lifetime budget constraint is the negative of the expression on RHS of \eqref{eq:SScost}.  Next revert to class notes for discussion of Ponzi Schemes.}}  

The reason the introduction of Social Security makes future 
generations worse off is that without SS they could have invested the 
amount $\Surplus_{1,t}$ and earned interest on it of $\rfree_{t+1}\Surplus_{1,t}$ in 
period $2$.  Now the money is taken away from them when young and 
returned {\it without interest} when old.  Thus, the loss is precisely 
the loss in interest income on $\Surplus_{1,t}$ in period $t+1$, discounted 
back to the present.

Note that if there is zero population growth, the foregoing analysis
all holds in per-capita terms as well, so that the per-capita
change in generational accounts from introducing Social Security 
is
\begin{equation}\begin{gathered}\begin{aligned}
        \bar{z}_{t} & =  \rfree_{t+1}\surplus_{1,t}/\Rfree_{t+1}
\end{aligned}\end{gathered}\end{equation}

\subsection{Effects of Population Growth}

If there is perpetual population growth, it is possible to finance a \opt{MarginNotes}{\marginpar{\tiny Note that the definition of $\surplus_{2,t+1}$ is slightly different from usual because we divide by $L_{t}$ rather than $L_{t+1}$.}}
positive rate of return on Social Security contributions.  Define
\begin{equation}
        \surplus_{1,t} = \Surplus_{1,t}/L_{t}
        \label{eq:zdef}
\end{equation}
and assume there is constant population growth, $\PopGro=L_{t+1}/L_{t}$.  
If we assume that Social Security taxes per capita are constant, 
then we can achieve a positive rate of return on Social Security 
contributions equal to the growth rate of population:
\begin{equation}\begin{gathered}\begin{aligned}
        \surplus_{2,t+1} & =  \Surplus_{2,t+1}/L_{t}  \\
                  & =  -\Surplus_{1,t+1}/L_{t}  \\
         & =  -\left(\frac{\Surplus_{1,t+1}}{L_{t+1}}\right)\left(\frac{L_{t+1}}{L_{t}}\right) 
\\       & =  -\surplus_{1,t+1}\PopGro = -\surplus_{1,t}\PopGro.
\end{aligned}\end{gathered}\end{equation}

Not only does this prove that it is {\it possible} for the Social 
Security system to pay a rate of return equal to the rate of 
population growth - it proves that the {\it only} rate of return that 
is consistent with constant per-capita taxes on the young is a rate of 
return of $\PopGro$.  \opt{MarginNotes}{\marginpar{\tiny It is not only possible, but {\it necessary} to pay a positive return in an economy with perpetual population growth.}}


\subsection{Effects of Productivity Growth and Population Growth}

Suppose there is wage growth ${\WGro}$ betwen $t$ and $t+1$, and 
suppose that workers contribute a constant {\it percentage} of their 
incomes to the Social Security system, $\surplus_{1,t} = \zeta \Wage_{1,t}$.  In 
this case it is possible to earn a rate of return on SS contributions 
equal to the product of the growth factor for wages and the growth 
factor for population:
\begin{equation}\begin{gathered}\begin{aligned}
        \surplus_{1,t} & =  \zeta \Wage_{1,t}  \\
        \Wage_{1,t+1} & =  {\WGro}\Wage_{1,t}  \\
        \surplus_{2,t+1} & =  -\Surplus_{1,t+1}/L_{t}
\\  & =  -(\Surplus_{1,t+1}/L_{t+1})(L_{t+1}/L_{t})
\\  & =  -\zeta \underbrace{\Wage_{1,t+1}}_{= {\WGro} \Wage_{1,t}} \PopGro
\\  & =  -\zeta \Wage_{1,t}{\WGro} \PopGro
\\  & =  -\surplus_{1,t} {\WGro}\PopGro
\end{aligned}\end{gathered}\end{equation}
so viewed from the perspective of the young generation in period $t$, 
their Social Security contributions are returned to them larger by a 
factor of ${\WGro}\PopGro$ than what they paid in; the effective rate of return is 
therefore ${\WGro}\PopGro$.

\subsection{Generational Accounts in a Growing Economy}

Now consider the per-capita generational accounts in an economy with 
constant population growth and constant wage growth and a Social 
Security system that imposes a constant tax of $\zeta$ on the wages of 
the young:
\begin{equation}\begin{gathered}\begin{aligned}
        \bar{z}_{t} & =  \surplus_{1,t}+\surplus_{2,t+1}/\Rfree_{t+1}
\\   & =  \zeta \Wage_{1,t}- {\WGro}\PopGro \zeta \Wage_{1,t}/\Rfree_{t+1}
\\   & =  \zeta \Wage_{1,t}\left(1 - {\WGro}\PopGro/\Rfree_{t+1}\right)
\\   & =  \zeta \Wage_{1,t}\left(\frac{\Rfree_{t+1} - {\WGro}\PopGro}{\Rfree_{t+1}}\right).
\end{aligned}\end{gathered}\end{equation}

Note that this expression will be {\it negative} if ${\WGro}\PopGro>\Rfree_{t+1}$, meaning 
that the introduction of a Social Security system with a positive tax 
rate $\zeta$ actually {\it improves} the lifetime budget constraint!  
This is another way of seeing that an economy is {\it dynamically 
inefficient} if the return factor for capital $\Rfree$ is less than the 
product of the population growth and productivity growth factors.  (Or, using approximations,
the rate of return is less than the sum of the population growth rate and the productivity growth rate).




  
\input handoutBibMake

\end{document}

It will often be convenient to write generation accounts in per-capita
terms rather than in aggregate terms.  Thus define
\begin{equation*}\begin{gathered}\begin{aligned}
        \surplus_{1,t} & =  \Surplus_{1,t}/L_{t}  \\
        \surplus_{2,t+1} & =  \Surplus_{2,t+1}/L_{t+1}.
\end{aligned}\end{gathered}\end{equation*}

